% from ega4-2-fr.tex
\oldpage[IV-2]{142}

(ii) Comme pour toute générisation $x'$ de $x\in f^{-1}(y)$,
$y'=f(x')$ est une générisation de $y$, on peut se borner à vérifier que si
$x\in\operatorname{Supp}(\mathscr G)$ et $f(x)=y$, on a
$\operatorname{prof}(\mathscr G_x)\geq\inf(k,\dim(\mathscr G_x))$; il résulte
de (6.1.2) et (6.3.1) que l'on a
\[
\dim(\mathscr G_x)=\dim(\mathscr E_y)+
\dim(\mathscr F_x\otimes_{\mathcal O_y}k(y))
\]
\[
\operatorname{prof}(\mathscr G_x)=\operatorname{prof}(\mathscr E_y)+
\operatorname{prof}(\mathscr F_x\otimes_{\mathcal O_y}k(y)).
\]

Par hypothèse, on a
\[
\operatorname{prof}(\mathscr E_y)\geq\inf(k,\dim(\mathscr E_y))
\]
\[
\operatorname{prof}(\mathscr F_x\otimes_{\mathcal O_y}k(y))
\geq\inf(k,\dim(\mathscr F_x\otimes_{\mathcal O_y}k(y)))
\]
d'où, en ajoutant membre à membre,
\begin{align*}
\operatorname{prof}(\mathscr G_x)
&\geq\inf(k,\dim(\mathscr E_y))
+\inf(k,\dim(\mathscr F_x\otimes_{\mathcal O_y}k(y)))\geq{}\\
&\geq\inf\bigl(k,\dim(\mathscr E_y)
+\dim(\mathscr F_x\otimes_{\mathcal O_y}k(y))\bigr)
=\inf(k,\dim(\mathscr G_x))
\end{align*}
ce qui prouve (ii).

\begin{corollary}[6.4.2]
\label{IV.2.s6.4.2-fr}
Soient $X$, $Y$ deux préschémas localement noethériens, $f:X\to Y$ un
morphisme plat, $\mathscr E$ un $\mathcal O_Y$-Module cohérent. Supposons que
pour tout $y\in Y$, le préschéma $f^{-1}(y)$ possède la propriété $(S_k)$.
Alors, pour que $f^*(\mathscr E)$ vérifie $(S_k)$ en un point $x$, il faut et
il suffit que $\mathscr E$ possède la propriété $(S_k)$ au point $f(x)$.
\end{corollary}

\begin{remark}[6.4.3]
\label{IV.2.s6.4.3-fr}
Soient $A$, $B$ deux anneaux locaux noethériens, $k$ le corps résiduel de
$A$, $\varphi:A\to B$ un homomorphisme local, $M$ un $B$-module de type fini
qui soit un $A$-module plat, et supposons en outre que le $A$-module $A$ et le
$(B\otimes_Ak)$-module $M\otimes_Ak$ possèdent la propriété $(S_n)$; peut-on
alors conclure que le $B$-module $M$ possède la propriété $(S_n)$? Nous
l'ignorons même lorsque $n=1$, $M=B$ et $B=\widehat A$, autrement dit nous
ignorons si en général le \emph{complété} d'un anneau local noethérien
vérifiant $(S_n)$ vérifie lui aussi $(S_n)$, même pour $n=1$. Posant
$Y=\operatorname{Spec}(A)$, $X=\operatorname{Spec}(B)$ et
$\mathscr F=\widetilde M$, il suffirait, en vertu de (6.4.1, (ii)), de montrer
que pour \emph{tout} $y\in Y$, $\mathscr F_y$ possède la propriété $(S_n)$
(et non seulement lorsque $y$ est le point fermé de $Y$); il suffirait aussi
de montrer que l'ensemble $U$ des $x\in X$ tels que
$\mathscr F_x\otimes_{\mathcal O_{f(x)}}k(f(x))$ vérifie $(S_k)$ est ouvert
dans $X$ (puisqu'il contient le point fermé de $X$ par hypothèse). Nous
démontrerons plus loin (12.1.4) que $U$ est ouvert quand $B$ est un anneau
local d'une $A$-algèbre de type fini et $M=B$. D'autre part, pour
$B=\widehat A$ et $M=B$, nous avons vu dans (6.3.8) que les préschémas
$f^{-1}(y)$ sont des préschémas de Cohen-Macaulay (autrement dit vérifient
$(S_n)$ pour \emph{tout} $n$) lorsque l'on suppose que $A$ est quotient d'un
anneau local de Cohen-Macaulay. On en conclut que lorsque $A$ est un quotient
d'un anneau local de Cohen-Macaulay (ou plus généralement d'un anneau local
vérifiant les hypothèses de (6.3.8)), \emph{pour que $A$ vérifie $(S_n)$, il
faut et il suffit que son complété $\widehat A$ vérifie $(S_n)$}. Il resterait
à voir si cette propriété subsiste sans restriction sur $A$.
\end{remark}

